\chapter{Implementation Details}
\label{chap:implementation}

This chapter details the concrete implementation of the autonomous lifecycle agent, including Pydantic Abstract Syntax Tree (\gls{ast}) data structures, the deterministic Lisp compiler, symbolic oracle wrappers, and agent orchestration.

\section{Pydantic AST Schemas (\texttt{DomainBlueprint})}
The intermediate representation is defined using Pydantic models in Python, enforcing strict structural type-checking and schema validation prior to compilation.

\subsection{Domain Blueprint Schema}
Listing~\ref{lst:pydantic_schema} illustrates the core Pydantic data structures defining domain predicates, action parameters, preconditions, and state-transition effects.

\begin{lstlisting}[language=Python, caption=Pydantic AST Data Structures for Domain Modeling, label=lst:pydantic_schema]
from typing import List, Dict, Optional
from pydantic import BaseModel, Field

class ParameterBlueprint(BaseModel):
    name: str = Field(description="Variable name with leading ?, e.g., ?v")
    type_name: str = Field(description="Target object type, e.g., truck")

class PredicateBlueprint(BaseModel):
    name: str
    parameters: List[ParameterBlueprint]

class ActionBlueprint(BaseModel):
    name: str
    parameters: List[ParameterBlueprint]
    precondition: List[str] = Field(description="Lisp S-expression preconditions")
    effect: List[str] = Field(description="Lisp S-expression state effects")

class DomainBlueprint(BaseModel):
    domain_name: str
    requirements: List[str] = Field(default_factory=lambda: [":strips", ":typing"])
    types: Dict[str, Optional[str]] = Field(description="Child to parent type dictionary")
    predicates: List[PredicateBlueprint]
    actions: List[ActionBlueprint]
\end{lstlisting}

\section{Deterministic Programmatic Lisp Compiler}
The programmatic compiler iterates over the validated \texttt{DomainBlueprint} instance and formats each section into valid PDDL S-expressions:

\begin{lstlisting}[language=Python, caption=Deterministic PDDL Compiler Snippet]
def compile_action(action: ActionBlueprint) -> str:
    params_str = " ".join([f"{p.name} - {p.type_name}" for p in action.parameters])
    precond_str = f"(and {' '.join(action.precondition)})" if len(action.precondition) > 1 else action.precondition[0]
    effect_str = f"(and {' '.join(action.effect)})" if len(action.effect) > 1 else action.effect[0]
    
    return (
        f"  (:action {action.name}\n"
        f"    :parameters ({params_str})\n"
        f"    :precondition {precond_str}\n"
        f"    :effect {effect_str})\n"
    )
\end{lstlisting}

Because string interpolation directly uses structured lists, parenthesis imbalance is physically impossible during compilation.

\section{Symbolic Oracle Wrappers}
The verification engine encapsulates external binary tools (\texttt{VAL} and \texttt{fast-downward}) via Python subprocess wrappers:

\begin{itemize}
    \item \textbf{VAL Integration:} Invokes \texttt{Validate -v domain.pddl problem.pddl plan.val}, capturing stdout for parsing state transition errors and predicate mismatch warnings.
    \item \textbf{Fast Downward Integration:} Executes \texttt{fast-downward --alias seq-sat-lama-2011 domain.pddl problem.pddl}, checking exit status codes and extracting search node expansion logs.
\end{itemize}

\section{Agent Orchestration via LangGraph}
The 4-phase lifecycle is orchestrated using LangGraph state graphs. The global agent state retains:
\begin{enumerate}
    \item \texttt{domain\_blueprint}: Current JSON \texttt{DomainBlueprint} AST.
    \item \texttt{compiled\_pddl}: Latest generated PDDL domain string.
    \item \texttt{verification\_status}: Oracle status (\texttt{VALID}, \texttt{SYNTAX\_ERROR}, \texttt{UNSOLVABLE}).
    \item \texttt{suspiciousness\_map}: Calculated fault scores $\sigma(n)$ across action nodes.
    \item \texttt{iteration\_count}: Feedback loop counter (max iterations = 5).
\end{enumerate}
